\chapter{Discussion, Conclusion, and Future Work}

\section{Interpretation of Results}
This final chapter interprets the experimental findings presented in Chapter 4 and connects them back to the research questions stated in Chapter 1. With all four planned experiments completed, the thesis can compare both architecture choice and training strategy within a single controlled pipeline. This makes it possible to discuss not only the value of pretrained representations, but also the situations in which scratch training may remain competitive.

Across the four experiments, ResNet50 with transfer learning achieved the best overall balance of predictive performance. It produced the highest test accuracy, recall, and F1-score, while DenseNet121 trained from scratch achieved the highest test AUC by a very small margin. DenseNet121 with transfer learning produced lower recall but the highest precision and specificity. This indicates that the architectures and training strategies reached different operating points under the same binary formulation. In practical terms, ResNet50 with transfer learning was better aligned with the screening-oriented objective of reducing missed abnormal cases, while DenseNet121 with transfer learning was stronger at avoiding unnecessary positive predictions.

These findings show that transfer learning is effective for this task, but not in a uniform way. For ResNet50, transfer learning produced clear gains over scratch training across accuracy, recall, F1-score, and AUC. For DenseNet121, however, the scratch model slightly outperformed the transfer-learning model in the main test metrics, while the transfer-learning model became more conservative and achieved higher precision and specificity. This reinforces the thesis claim that architecture choice and training strategy should be evaluated together rather than treated as independent design choices.

\section{Implications for Medical Image Classification}
The results have several implications for practical chest X-ray classification research at the capstone level. First, they show that a relatively focused binary classification problem can still surface meaningful tradeoffs between model architectures. Even when two models achieve similar AUC values, their recall, precision, and specificity can differ substantially. This is important because medical imaging tasks are rarely judged by accuracy alone. In a clinical screening setting, a model that misses too many abnormal images may be less useful even if its overall accuracy remains acceptable.

Second, the Grad-CAM results provide moderate evidence that the selected best-balance model learned clinically plausible spatial patterns. In the correctly classified cases, the heatmaps concentrated mainly on thoracic regions instead of obvious background areas, which supports the claim that the model relied on image content relevant to the chest. This improves the interpretability of the result and makes the model behavior easier to explain in the thesis and in the defense. However, the false positive and false negative examples also demonstrate that plausible visual attention does not guarantee correct reasoning. Grad-CAM should therefore be treated as a transparency tool rather than as proof that the model has learned clinically reliable causal structure.

Third, the completed experiments reinforce the value of keeping the project scope controlled. By restricting the task to chest X-ray binary classification with two architectures and one explainability method, the project was able to produce complete training results, interpretable visual evidence, and a clear comparison that can be defended academically. This is preferable to attempting a much broader medical imaging project that would risk incomplete implementation or weak experimental depth.

\section{Limitations}
Several limitations should be acknowledged when interpreting the final results. The first is dataset quality. NIH ChestXray14 is a widely used public benchmark, but its labels were derived automatically from radiology reports rather than assigned manually for this exact classification task. As a result, label noise is likely present. This limits the confidence with which the measured performance can be interpreted as a true estimate of clinical diagnostic quality.

The second limitation is the binary task formulation itself. In this thesis, all non-``No Finding'' labels were grouped into a single abnormal class. This simplification was useful for making the capstone achievable and for supporting a clean comparison across models, but it also removes important clinical detail. Different abnormalities have different visual signatures and different levels of difficulty. A model that performs reasonably on binary normal-versus-abnormal classification may still struggle on more realistic multi-label or disease-specific tasks.

The third limitation is experimental coverage. Although the full four-way comparison was completed, the project still examines only two model architectures and one binary task formulation. The conclusions therefore should not be generalized to all deep learning approaches for chest X-ray analysis. Additional architectures, alternative thresholds, and disease-specific tasks could reveal different performance patterns.

The fourth limitation concerns explainability. Grad-CAM provides an intuitive visualization of the image regions that most influence the prediction, but it remains a qualitative method. Heatmaps can look plausible even when the underlying prediction is wrong, and different explainability methods may highlight different regions. For this reason, the Grad-CAM analysis in this thesis should be interpreted as supportive evidence about model attention rather than as a definitive account of how the model reasons.

\section{Conclusion}
This thesis investigated whether transfer learning and explainable deep learning could improve chest X-ray diagnostic classification in a focused capstone setting. The study used the NIH ChestXray14 dataset and formulated the task as binary classification, where images labeled ``No Finding'' were treated as normal and images with any pathology label were treated as abnormal. Two convolutional neural network architectures, ResNet50 and DenseNet121, were selected as the main comparison models because both are well established in image classification and have been widely used in medical imaging research. Four controlled experiments were completed in total: ResNet50 from scratch, ResNet50 with transfer learning, DenseNet121 from scratch, and DenseNet121 with transfer learning.

The final experimental results show that the best outcome depends on how performance is defined. DenseNet121 trained from scratch achieved the highest test AUC at 0.7457, while ResNet50 with transfer learning delivered the strongest overall balance of accuracy, recall, and F1-score, reaching a test accuracy of 0.6955, F1-score of 0.6547, and AUC of 0.7454. DenseNet121 with transfer learning produced a similar AUC of 0.7421 and the highest precision and specificity, but at the cost of noticeably lower recall. ResNet50 from scratch served as the weakest of the four runs, with a test AUC of 0.7225. These results indicate that the four configurations reached different tradeoff points under the same binary task formulation, and that transfer learning was clearly beneficial for ResNet50 but not uniformly beneficial for DenseNet121.

The Grad-CAM analysis further supported the practical value of the selected best-balance model. In representative correct predictions, the heatmaps concentrated primarily on thoracic regions rather than on irrelevant image background, suggesting that the model learned clinically plausible visual patterns. At the same time, the false positive and false negative cases showed that plausible attention maps do not guarantee correct decisions. This confirms that explainability should be used to support interpretation of model behavior rather than to replace quantitative evaluation.

Overall, the main implication of this thesis is that a carefully scoped workflow combining controlled experimentation, reproducible evaluation, and explainability can produce academically meaningful results for medical image analysis at the capstone level. Even with a simplified binary problem and only two model architectures, the study was able to reveal clear tradeoffs between sensitivity, specificity, precision, architecture choice, and training strategy. Most importantly, the results show that transfer learning improved ResNet50 substantially, but its benefit was not uniform across architectures. Grad-CAM also improved interpretability by helping connect quantitative model behavior with representative correct and incorrect predictions.

\section{Future Work}
Several directions can extend this work in a more advanced study. First, the binary classification setting can be expanded into multi-label classification so that individual thoracic diseases are modeled directly instead of being merged into a single abnormal class. This would make the task more clinically realistic and allow model performance to be analyzed at the disease level.

Second, the experimental comparison can be broadened. Additional architectures such as EfficientNet, vision transformers, or chest X-ray models trained with self-supervised learning could be evaluated to determine whether the observed ResNet50--DenseNet121 tradeoff generalizes to more recent model families. Future work should also test separate learning-rate schedules for scratch training and transfer-learning fine-tuning, since pretrained models may benefit from a more conservative fine-tuning rate than randomly initialized models. Threshold optimization and calibration should also be explored so that the present accuracy-recall-specificity tradeoffs can be tuned more explicitly for screening or decision-support use cases.

Third, explainability can be improved beyond the present Grad-CAM analysis. Future work could compare Grad-CAM with other explainability methods, incorporate uncertainty estimation, or involve expert review of model heatmaps. This would help determine whether the highlighted regions are not only visually plausible, but also clinically meaningful from a radiological perspective.

Finally, the generalizability of the findings should be tested on additional datasets and settings. External validation on benchmarks such as CheXpert or other institutional chest X-ray collections would provide stronger evidence about robustness across domains. If the long-term goal is clinical deployment or decision support, future work should also examine threshold selection, calibration, domain shift, and user-facing evaluation rather than relying on a single fixed test setting.
